%% kapa: exam=1 round=37 year=2026 period=1 form=A subject=civil_law subject_ko=민법 seq=7 number=7
%% source: Q-Net 감정평가사 기출문제 (cst00309) artlSeq=5258924 file=raw/1008/5258924/2026년 제37회 감정평가사 1차 1교시.pdf
%% conversion: kordoc 4.24.51 -> tools/gen1.py (auto); review=proofread-sonnet; flags=none
\begin{question}{7}
비법인사단 甲과 그 대표자 乙에 관한 설명으로 옳은 것을 모두 고른 것은? (다툼이 있으면 판례에 따름)
\begin{framedtext}
ㄱ. 甲이 총유물에 관한 매매계약을 체결한 경우, 乙이 그 계약에 따라 발생한 채무에 대해 소멸시효 중단의 효력이 있는 승인을 하는 것은 총유물의 관리ㆍ처분행위이다. \par
ㄴ. 乙이 정관에서 정한 대표권 제한을 위반하여 丙과 거래행위를 한 경우, 丙이 그 대표권 제한사실을 알았거나 알 수 있었을 경우가 아니라면 그 거래행위는 유효하다. \par
ㄷ. 乙로부터 대표자로서의 모든 권한을 포괄적으로 위임받은 丁이 甲을 대표하여 한 대행행위의 효력은 원칙적으로 甲에게 미치지 않는다. \par
\end{framedtext}
\begin{choices}
  \item ㄱ
  \item ㄷ
  \item ㄱ, ㄴ
  \item ㄴ, ㄷ
  \item ㄱ, ㄴ, ㄷ
\end{choices}
\end{question}
\answer{4}
